% Lecture 1 — Complex numbers
% Compiled to web slides by tools/tex2slides.mjs

\lecture{Complex numbers}{Linear Algebra · Warsaw University of Technology}{Lecture 1 · Tomasz Brengos}

\section{Motivation}

\begin{frame}{A problem with no real solution}
The equation
\[
x^2 + 1 = 0
\]
has no solution in $\R$, since $x^2 \geq 0$ for every real $x$.

\pause

The remedy: \emph{extend} the real numbers by one new number
\[
i, \qquad \text{required to satisfy} \qquad i^2 = -1.
\]

\pause

\begin{remark}
Adjoining $i$ to $\R$ produces the field $\C$ — in which \emph{every} polynomial
equation has a solution. That is where this lecture ends.
\end{remark}
\end{frame}

\begin{frame}{Born of a scandal: the cubic (Italy, 1545)}
Historically, nobody extended $\R$ because of $x^2+1=0$ — a quadratic with no
solution was simply \emph{unsolvable}, end of story.

\pause

The trouble began with \emph{cubics}. Del Ferro and Tartaglia found a solution
formula, published by Cardano in his \emph{Ars Magna} (1545): the equation
$x^3 = px + q$ is solved by
\[
x = \sqrt[3]{\tfrac{q}{2} + \sqrt{\tfrac{q^2}{4} - \tfrac{p^3}{27}}}
  + \sqrt[3]{\tfrac{q}{2} - \sqrt{\tfrac{q^2}{4} - \tfrac{p^3}{27}}}.
\]

\pause

\begin{warning}
Try it on $x^3 = 15x + 4$. The formula demands
$x = \sqrt[3]{2 + \sqrt{-121}} + \sqrt[3]{2 - \sqrt{-121}}$ — a square root of
a \emph{negative} number. Yet $x = 4$ plainly solves the equation:
$64 = 60 + 4$. The formula is right, and it is speaking a language nobody
understood yet.
\end{warning}
\note{Key historical point: complex numbers were forced on us by cubics with perfectly real roots — not by quadratics.}
\end{frame}

\begin{frame}{Bombelli's wild thought (1572)}
Rafael Bombelli had no theory of complex numbers — no definition, no $i$.
His idea (``a wild thought'', he admitted): treat the \emph{meaningless
symbol} $\sqrt{-1}$ as if it were a number, obey every ordinary rule of
algebra, and add one bookkeeping rule: $(\sqrt{-1})^2 = -1$.

\pause

\begin{steps}
\step[expand the cube by the binomial rule — pure symbol pushing]{\left(2+\sqrt{-1}\right)^3 = 8 + 12\sqrt{-1} + 6\left(\sqrt{-1}\right)^2 + \left(\sqrt{-1}\right)^3}
\step[apply the bookkeeping rule twice]{= 8 + 12\sqrt{-1} - 6 - \sqrt{-1} = 2 + 11\sqrt{-1} = 2 + \sqrt{-121}}
\step[the same for the other cube root]{\sqrt[3]{2 + \sqrt{-121}} = 2 + \sqrt{-1}, \qquad \sqrt[3]{2 - \sqrt{-121}} = 2 - \sqrt{-1}}
\step[the symbols cancel — and the answer is checkable]{x = \left(2 + \sqrt{-1}\right) + \left(2 - \sqrt{-1}\right) = 4}
\end{steps}

\pause

\begin{remark}
The symbol entered the computation, cancelled out, and delivered the honest,
verifiable answer $4$ — with no claim about what $\sqrt{-1}$ \emph{is}.
That claim is a debt, and this lecture pays it: our definition of $\C$ is
precisely the licence for Bombelli's game.
\end{remark}
\end{frame}

\begin{frame}{Two centuries of suspicion}
\begin{itemize}
\item \textbf{1637} — Descartes dismisses the new numbers as \emph{imaginary};
the insult sticks and becomes the official name.
\pause
\item \textbf{1748} — Euler connects them to trigonometry:
$e^{i\varphi} = \cos\varphi + i\sin\varphi$; in \textbf{1777} he introduces the
symbol $i$.
\pause
\item \textbf{1797 / 1806} — Wessel and Argand independently draw complex
numbers as \emph{points of the plane}.
\pause
\item \textbf{1831} — Gauss champions the planar picture, coins the name
\emph{complex numbers}, and the suspicion finally dies.
\pause
\item \textbf{1833} — Hamilton gives the modern definition: pairs $(a,b)$ of
real numbers. Ten years later his quaternions — a 4-dimensional cousin of
$\C$ — end up rotating every 3D game world you have played in.
\end{itemize}
\pause
Today's lecture follows the same arc: calculate first, then draw the plane —
in one lecture instead of three hundred years.
\end{frame}

\begin{frame}{You have used complex numbers today}
\begin{itemize}
\item \textbf{Multimedia.} Every JPEG, MP3 and video call runs on the Fast Fourier
Transform — arithmetic of the \emph{roots of unity} we compute today.
\pause
\item \textbf{Graphics and games.} Multiplying by a complex number of modulus $1$
\emph{rotates} the plane; one dimension up, quaternions rotate every 3D game world.
\pause
\item \textbf{Quantum computing.} A qubit is a pair of complex numbers; quantum
interference is addition of complex numbers with different arguments.
\pause
\item \textbf{Fractals.} The Mandelbrot set is the iteration $z \mapsto z^2 + c$,
coloured by how fast $|z| \to \infty$.
\end{itemize}
\note{Sell the chapter: FFT and rotations are the two examples to come back to later.}
\end{frame}

\section{The complex numbers}

\begin{frame}{Definition}
\begin{definition}
A \emph{complex number} is an expression $a + bi$, where $a, b \in \R$.
The set of all complex numbers is denoted by $\C$.
\end{definition}

\pause

\begin{remark}
Formally a complex number is just the pair $(a,b) \in \R \times \R$: the symbol $i$
denotes $(0,1)$, and nothing is being ``imagined'' — $\C$ is the plane $\R^2$
equipped with a multiplication.
\end{remark}
\end{frame}

\begin{frame}{Arithmetic}
\begin{definition}
\[
(a+bi) + (c+di) = (a+c) + (b+d)i, \qquad
(a+bi)(c+di) = (ac - bd) + (ad + bc)i.
\]
\end{definition}

The product rule is what formal expansion gives, using $i^2 = -1$:

\pause

\begin{steps}
\step[expand formally]{(2+3i)(1-2i) = 2 - 4i + 3i - 6i^2}
\step[use i^2 = -1]{= 2 - i + 6}
\step[collect]{= 8 - i}
\end{steps}
\end{frame}

\begin{frame}{Division}
\begin{fact}
For any complex number, $(a+bi)(a-bi) = a^2 + b^2$ — a real number.
Hence every $z \neq 0$ has a multiplicative inverse.
\end{fact}

\pause

\begin{steps}
\step[multiply by the conjugate]{\frac{1}{1+i} = \frac{1-i}{(1+i)(1-i)}}
\step[the denominator became real]{= \frac{1-i}{2} = \tfrac12 - \tfrac12\, i}
\end{steps}

\pause

\begin{theorem}
$\C$ with these operations is a \emph{field}: all rules of arithmetic
familiar from $\Q$ and $\R$ hold. (Fields are the subject of the next lecture.)
\end{theorem}
\end{frame}

\section{The complex plane}

\begin{frame}{A number is a point}
\begin{columns}
\column
A complex number $z = a + bi$ \emph{is} a point of the plane.

Its length is the \emph{modulus}
\[
|z| = \sqrt{a^2 + b^2},
\]
and the angle to the positive real axis is the \emph{argument} $\arg(z)$.

\textbf{Drag the point.}
\column
\widget{CPlane}{re=2, im=1.5}
\end{columns}
\end{frame}

\begin{frame}{Conjugation is reflection}
\begin{columns}
\column
\begin{definition}
The \emph{conjugate} of $z = a+bi$ is $\bar z = a - bi$.
\end{definition}

\begin{fact}
$z \bar z = |z|^2$, \quad $\Re(z) = \tfrac12(z + \bar z)$, \quad
$\overline{z+w} = \bar z + \bar w$, \quad $\overline{zw} = \bar z\, \bar w$.
\end{fact}

Geometrically: reflection across the real axis.
\column
\widget{CPlane}{re=2, im=1.5, showConj=true, showPolar=false}
\end{columns}
\end{frame}

\begin{frame}{Moduli measure distance}
\begin{fact}
$|z - w|$ is the distance between the points $z$ and $w$, and the
\emph{triangle inequality} holds: $|z+w| \leq |z| + |w|$.
\end{fact}

\pause

\begin{columns}
\column
Conditions on moduli describe geometric figures:
\[
\{\, z \in \C \mid |z - z_0| = r \,\}
\]
is the \emph{circle} of radius $r$ centred at $z_0$.

For $z_0 = 0$ this is $|z| = r$, centred at the origin. A general $z_0$
shifts the centre — mind the signs: $|z + 1 - i| = |z - (-1+i)|$ is centred
at $-1+i$, \emph{not} at $1-i$.

Drag $z$ onto the dashed circle: the distance $|z - z_0|$ hits $1.5$.
\column
\widget{CPlane}{re=0.5, im=1, circleRe=-1, circleIm=1, circleR=1.5, showPolar=false}
\end{columns}
\end{frame}

\section{Polar form}

\begin{frame}{The argument, precisely}
\begin{definition}
For $z \neq 0$, an \emph{argument} of $z$ is any angle $\varphi$ with
\[
\cos\varphi = \frac{a}{|z|}, \qquad \sin\varphi = \frac{b}{|z|}.
\]
The argument in $(-\pi, \pi]$ is the \emph{principal} argument.
By convention $\arg(0) = 0$.
\end{definition}

\pause

\begin{warning}
$\tan\varphi = \tfrac{b}{a}$ alone does \emph{not} determine $\varphi$: always
check the quadrant. For $z = -1 + i$ the argument is $\tfrac{3\pi}{4}$,
not $-\tfrac{\pi}{4}$.
\end{warning}
\end{frame}

\begin{frame}{Polar form}
\begin{theorem}
Every $z \in \C$ can be written as
\[
z = |z| \left( \cos\varphi + i \sin\varphi \right), \qquad \varphi = \arg(z).
\]
For $z = 0$ both sides are $0$ whatever the angle — which is why the
convention $\arg(0) = 0$ costs nothing.
\end{theorem}

\pause

\begin{example}
$z = 1 + i$: modulus $|z| = \sqrt2$, argument $\tfrac{\pi}{4}$, so
\[
1 + i = \sqrt2 \left( \cos\tfrac{\pi}{4} + i \sin\tfrac{\pi}{4} \right).
\]
\end{example}
\end{frame}

\begin{frame}{Multiplication is rotation and scaling}
\begin{theorem}
$z_1 z_2 = |z_1||z_2| \big( \cos(\varphi_1 + \varphi_2) + i \sin(\varphi_1 + \varphi_2) \big)$
— \emph{moduli multiply, arguments add}.
\end{theorem}

\widget{ComplexMul}{}
\note{Let students drag w onto the unit circle: pure rotation, the heart of computer graphics.}
\end{frame}

\begin{frame}{De Moivre's formula}
\begin{theorem}
If $z = |z|(\cos\varphi + i\sin\varphi)$ and $n \in \N$, then
$z^n = |z|^n \big( \cos(n\varphi) + i \sin(n\varphi) \big)$.
\end{theorem}

\pause

\begin{steps}
\step[polar form of 1+i]{(1+i)^{100} = \left(\sqrt2\right)^{100} \left( \cos\tfrac{100\pi}{4} + i \sin\tfrac{100\pi}{4} \right)}
\step[De Moivre]{= 2^{50} \left( \cos 25\pi + i \sin 25\pi \right)}
\step[angles live modulo 2 pi]{= 2^{50} (\cos\pi + i\sin\pi) = -2^{50}}
\end{steps}
\end{frame}

\section{Roots}

\begin{frame}{Every number has n distinct n-th roots}
\begin{theorem}
For $z = |z|(\cos\varphi + i\sin\varphi) \neq 0$, the solutions of $w^n = z$ are
\[
w_k = \sqrt[n]{|z|} \left( \cos\tfrac{\varphi + 2k\pi}{n} + i \sin\tfrac{\varphi + 2k\pi}{n} \right),
\qquad k = 0, 1, \ldots, n-1.
\]
\end{theorem}

\pause

\begin{fact}
The $n$ roots lie on a circle of radius $\sqrt[n]{|z|}$ and cut it into $n$
equal arcs — a regular $n$-gon.
\end{fact}
\end{frame}

\begin{frame}{Roots of unity}
\begin{columns}
\column
Taking $z = 1$: the $n$-th \emph{roots of unity}
\[
z_k = \cos\tfrac{2k\pi}{n} + i \sin\tfrac{2k\pi}{n}.
\]

Multiplying by $z_1$ rotates the polygon by one step — press the button.

These numbers drive the Fast Fourier Transform.
\column
\widget{RootsOfUnity}{n=6}
\end{columns}
\end{frame}

\begin{frame}{Example: the square roots of 2i}
\begin{steps}
\step[write 2i in polar form]{2i = 2 \left( \cos\tfrac{\pi}{2} + i \sin\tfrac{\pi}{2} \right)}
\step[k = 0]{w_0 = \sqrt2 \left( \cos\tfrac{\pi}{4} + i \sin\tfrac{\pi}{4} \right) = 1 + i}
\step[k = 1]{w_1 = \sqrt2 \left( \cos\tfrac{5\pi}{4} + i \sin\tfrac{5\pi}{4} \right) = -1 - i}
\end{steps}

\pause

Check: $(1+i)^2 = 1 + 2i + i^2 = 2i$. ✓
\end{frame}

\section{Polynomials}

\begin{frame}{The Fundamental Theorem of Algebra}
\begin{theorem}
Every polynomial $f \in \C[x]$ with $\deg(f) \geq 1$ has a root in $\C$.
\end{theorem}

\pause

\begin{corollary}
Every such $f$ factors completely into linear factors:
\[
f(x) = a (x - z_1)^{k_1} \cdots (x - z_m)^{k_m}.
\]
\end{corollary}

\pause

This is why we extended $\R$ in the first place: over $\C$, polynomial
equations always have solutions — the extension is \emph{complete}.
\end{frame}

\begin{frame}{Quadratic equations over C}
\begin{theorem}
For $f(x) = ax^2 + bx + c$ with $a \neq 0$, the roots are
$z_{0,1} = \frac{-b \mp \delta}{2a}$, where $\delta$ is any square root of
$\Delta = b^2 - 4ac$.
\end{theorem}

\pause

\begin{steps}
\step[negative discriminant is no obstacle]{x^2 + 2x + 5: \quad \Delta = 4 - 20 = -16, \qquad \delta = 4i}
\step[conjugate pair]{z_{0,1} = \frac{-2 \mp 4i}{2} = -1 \mp 2i}
\end{steps}
\end{frame}

\begin{frame}{Real polynomials and conjugate roots}
\begin{theorem}
If $f \in \R[x]$ and $f(z) = 0$, then also $f(\bar z) = 0$: non-real roots of
real polynomials come in conjugate pairs.
\end{theorem}

\pause

\begin{corollary}
Every real polynomial factors into real linear factors and real quadratic
factors with $\Delta < 0$.
\end{corollary}

\pause

\begin{example}
$x^4 - 16 = (x-2)(x+2)(x^2+4)$ over $\R$; over $\C$ the last factor splits
into $(x-2i)(x+2i)$ — a conjugate pair.
\end{example}
\end{frame}

\begin{frame}{Summary}
\begin{itemize}
\item $\C$ = the plane with a multiplication; $i^2 = -1$; conjugate and modulus.
\pause
\item Polar form: multiplication \emph{rotates and scales}; De Moivre computes powers.
\pause
\item Every $z \neq 0$ has exactly $n$ $n$-th roots — a regular $n$-gon.
\pause
\item Fundamental Theorem of Algebra: over $\C$, polynomials factor completely.
\pause
\item \textbf{Next lecture:} what exactly is a \emph{field}? Groups, fields, and
arithmetic modulo $p$.
\end{itemize}
\end{frame}
